% ============================================================================
% Chapter 29: Working with Worksheets
% ============================================================================
\chapter{Working with Worksheets}
\label{ch:excel-worksheets}

\begin{objectives}
  \item Add, rename, delete, and rearrange worksheets
  \item Enter different types of data (text, numbers, dates)
  \item Adjust column widths and row heights
  \item Use AutoFill to quickly enter data series
\end{objectives}

\bytetip[tip]{A workbook is like a notebook with multiple pages. Each page is a \textbf{worksheet} (or ``sheet''). You can have as many sheets as you need!}

\section{Managing Worksheets}
\label{sec:manage-sheets}
\index{MS Excel!worksheets}

\begin{theorybox}[Worksheet Operations]
\begin{tabular}{ll}
  \textbf{Action} & \textbf{How To Do It} \\
  \midrule
  Add a new sheet & Click the \textbf{+} button next to sheet tabs \\
  Rename a sheet & Double-click the sheet tab and type a new name \\
  Delete a sheet & Right-click the tab $\rightarrow$ Delete \\
  Move a sheet & Click and drag the tab to a new position \\
  Copy a sheet & Right-click $\rightarrow$ Move or Copy $\rightarrow$ check ``Create a copy'' \\
\end{tabular}
\end{theorybox}

\section{Entering Data}
\label{sec:entering-data}
\index{MS Excel!data entry}

Excel cells can hold three types of data:
\begin{itemize}[leftmargin=*, label={\textcolor{ModuleBlue}{\faAngleRight}}, itemsep=3pt]
  \item \textbf{Text} (labels) --- aligned to the left by default (e.g., ``Name'', ``Total'')
  \item \textbf{Numbers} --- aligned to the right by default (e.g., 42, 3.14, -5)
  \item \textbf{Dates} --- displayed in date format (e.g., 13/04/2026)
\end{itemize}

\section{AutoFill}
\label{sec:autofill}
\index{MS Excel!AutoFill}

\hl{AutoFill} is one of Excel's most powerful features. It automatically fills cells with a pattern you start.

\begin{labstep}[Using AutoFill]
\begin{enumerate}[label=\protect\stepcircle{\arabic*}, leftmargin=2cm, itemsep=3pt]
  \item Type \textbf{1} in cell A1 and \textbf{2} in cell A2.
  \item Select both cells (A1:A2).
  \item Hover over the small square at the bottom-right corner (the \textbf{fill handle}).
  \item Click and drag down to A10.
  \item Excel automatically fills: 3, 4, 5, 6, 7, 8, 9, 10!
\end{enumerate}
\end{labstep}

AutoFill also works with: days of the week, months, dates, and custom patterns.

\bytetip[fun]{AutoFill is smart enough to continue patterns like ``Mon, Tue, Wed\ldots'' or ``Jan, Feb, Mar\ldots'' It even works with ``Quarter 1, Quarter 2, Quarter 3\ldots'' --- try it!}

\begin{funfact}
Excel's grid has been the same size since Excel 2007: 16,384 columns and 1,048,576 rows. Before that (Excel 2003), it only had 256 columns and 65,536 rows. The upgrade was a big deal for scientists and data analysts!
\end{funfact}

\begin{reviewqs}
  \item How do you rename a worksheet?
  \item What three types of data can you enter in a cell?
  \item What is AutoFill and how do you use it?
  \item How do you add a new worksheet to a workbook?
  \item How can you adjust the width of a column?
\end{reviewqs}

\begin{challenge}
Create a workbook with 3 sheets: ``Names'', ``Marks'', and ``Summary.'' In ``Names,'' use AutoFill to create a numbered list (1--20) in column A. In column B, add 20 names. Rename all three sheet tabs with meaningful names.
\end{challenge}
